\documentclass[../DoAn.tex]{subfiles}
\begin{document}
Chương này có độ dài không quá 10 trang. Chương này sinh viên trình bày về phân tích rõ ngữ cảnh bài toán cũng như các kết quả nghiên cứu tương tự. Đồng thời, sinh viên có thể trình bày thêm về các kiến thức nền tảng. Mỗi chương nên có đoạn mở đầu chươnggiới thiệu những nội dung sẽ trình bày trong chương.

\section{Ngữ cảnh của bài toán}
\ldots

\section{Các kết quả nghiên cứu tương tự}
Trong phần này sinh viên trình bày các nghiên cứu liên quan (related work), chú ý phân tích rõ những ưu nhược điểm của chúng. Từ đó, nêu bật lên động lực để thực hiện nghiên cứu của đồ án này.

\section{Linux Networking Stack}

\section{eBPF}
\subsection{Lịch sử và tổng quan}
\subsection{Cơ chế hoạt động}
\subsection{BPF Maps}

\section{XDP}
\subsection{Tổng quan và vị trí trong kernel}
\subsection{Các chế độ hoạt động}
\subsection{Các return code}
\subsection{So sánh hiệu năng với iptables}

\section{Firewall phân tầng}
\subsection{Lọc gói tin tầng 3, tầng 4}
\subsection{Lọc trạng thái tầng 4}
\subsection{Kiểm tra gói tin tầng 7}

\section{Các dạng tấn công từ chối dịch vụ}
\subsection{Tấn công SYN Flood}
\subsection{Tấn công UDP Flood}
\subsection{Tấn công HTTP Flood}
\subsection{Slowloris}
\subsection{IP Spoofing và Random Source Flooding}

\section{Random Early Detection (RED) và Weighted RED (WRED)}
\subsection{Cơ chế Random Early Detection (RED)}
\subsection{Weighted RED và ưu tiên theo loại traffic}
\subsection{Ứng dụng Probabilistic Drop trong môi trường DDoS}

\section{GeoIP, ASN và định dạng MMDB}
\subsection{Tên của kiến thức nền tảng số 2.1}

\section{Iptables và Nftables}

\section{Suricata}

\section{ Connection Tracking (Conntrack)}


\textbf{Lưu ý}: Mỗi chương nên có đoạn kết thúc chương,  tổng kết lại các nội dung đã trình bày.

\end{document}